\clearpage
\begin{flushright}
  \fechaclase{11 de septiembre de 2026}
\end{flushright}

Considere que
\[
  \dot{V}\leq-\alpha V^{1/2},
  \qquad \alpha>0
\]
para algún $\alpha$ arbitrario, despejando
\[
  V^{-1/2}\dot{V}\leq-\alpha
\]
\[
  \int_{V(0)}^{V(t)}V^{-1/2}\,dV
  \leq-\int_0^t\alpha\,dt
\]
\[
  2V^{1/2}-2V^{1/2}(0)\leq-\alpha t
\]
como $V\to0$
\[
  t_r=\frac{2V^{1/2}(0)}{\alpha}
  =\frac{\lvert\sigma(0)\rvert}{\alpha/\sqrt{2}}
\]

\begin{center}
\begin{tikzpicture}[scale=0.78,>=Stealth]
  % Ejes y superficie de deslizamiento.
  \draw[->,very thick] (-3.4,0)--(3.35,0) node[right] {$X_1$};
  \draw[->,very thick] (0,-3.25)--(0,3.25) node[above] {$X_2$};
  \draw[green!60!black,very thick,dashed] (-1.95,3.45)--(2.1,-3.45);

  % Origen.
  \fill[orange] (0,0) circle (3pt);
  \node[orange,above right] at (0.05,0.05) {$(0,0)$};

  % Trayectoria superior: punto inicial -> superficie.
  \coordinate (Xu) at (1.35,2.25);
  \coordinate (Su) at (-1.22,2.0);
  \fill[orange] (Xu) circle (4pt);
  \node[orange,right] at (1.48,2.15) {$X(CX_{10},X_2)$};
  \draw[rojotrayectoria,very thick,
    postaction={decorate},
    decoration={markings,
      mark=at position 0.28 with {\arrow{Stealth}},
      mark=at position 0.68 with {\arrow{Stealth}}}]
    (Xu) .. controls (0.75,3.0) and (-0.55,3.0) .. (Su);
  \node[orange,left] at (-1.25,2.03) {$\sigma=0$};

  % Movimiento deslizante superior -> origen.
  \draw[rojotrayectoria,very thick,decorate,-{Stealth},
    decoration={zigzag,segment length=4pt,amplitude=2.4pt}]
    (Su)--(0,0);

  % Trayectoria inferior: punto inicial -> superficie.
  \coordinate (Xl) at (-0.72,-3.0);
  \coordinate (Sl) at (1.65,-2.7);
  \fill[orange] (Xl) circle (4pt);
  \node[orange,above left] at (-0.72,-2.92) {$X$};
  \draw[moradotrayectoria,very thick,
    postaction={decorate},
    decoration={markings,
      mark=at position 0.25 with {\arrow{Stealth}},
      mark=at position 0.55 with {\arrow{Stealth}},
      mark=at position 0.82 with {\arrow{Stealth}}}]
    (Xl) .. controls (-0.15,-3.65) and (1.2,-3.65) .. (Sl);

  % Movimiento deslizante inferior -> origen.
  \draw[moradotrayectoria,very thick,decorate,-{Stealth},
    decoration={zigzag,segment length=4pt,amplitude=2.4pt}]
    (Sl)--(0,0);

  % Explicación a la derecha.
  \node[blue,align=left,anchor=west,text width=6.8cm] at (4.8,1.0)
    {Lo que se busca es que del\\
     punto inicial $(X)$ mandarlo a la\\
     superficie $(\sigma)$, y de la superficie\\
     lo manda a $(0,0)$};
  \node[rojotrayectoria,anchor=west] at (4.8,-0.85) {Trayectoria 1};
  \node[moradotrayectoria,anchor=west] at (4.8,-1.3) {Trayectoria 2};
\end{tikzpicture}
\end{center}

\[
  \sigma=X_2+cX_1=0
\]
\[
  y+mx=0
\]
\[
  y=-mx
\]
\[
  \sigma=X_2=-cX_1
\]

\paragraph{\tituloejemplo{Ejemplo 1 SMC}}
Diseñe un controlador SMC para el sistema
\[
  \dot{X}_1=X_2
\]
\[
  \dot{X}_2=
  \underbrace{\operatorname{sen}(X_1+X_2)+\cos(X_1+X_2)}_{f(X)}+v
\]
\[
  \sigma=X_2+cX_1
\]
\[
  v=-cX_2-K\operatorname{sign}(\sigma)
\]

$K$ debe cumplir $K>L$, recordemos que
\[
  \lvert f(X,t)\rvert\leq L
\]
por lo que
\[
  K>2
\]